\chapter{Classical Mechanics}\label{classical-mechanics}
\chapterstatus{planned}

Classical mechanics is where much of analysis and geometry was born: the calculus of variations,
symplectic geometry and dynamical systems all began as ways of solving Newton's equations. This chapter
presents mechanics as an application of those theories, from Newton's laws to integrable systems.

\noindent\textbf{Prerequisites.} Chapter~\ref{ordinary-differential-equations} (Ordinary Differential Equations and Dynamical Systems), Chapter~\ref{calculus-of-variations} (Calculus of Variations), Chapter~\ref{smooth-manifolds} (Smooth Manifolds), Chapter~\ref{lie-groups} (Lie Groups), Chapter~\ref{riemannian-geometry} (Riemannian Geometry) and Chapter~\ref{symplectic-geometry} (Symplectic Geometry).

\noindent\textbf{Planned contents.}
\begin{enumerate}
\item Newtonian mechanics: Galilean spacetime, conservative forces, energy and angular momentum; the
harmonic oscillator; the Kepler problem, Kepler's laws and the Runge--Lenz vector; Bertrand's theorem.
\item Lagrangian mechanics on configuration manifolds: constraints, the principle of least action,
mechanical systems as geodesic flows (Jacobi's metric); Noether's theorem and conservation laws.
\item Small oscillations and normal modes; stability of equilibria (Lagrange--Dirichlet).
\item Rigid bodies: the configuration space $\mathrm{SO}(3)$, the inertia tensor, Euler's equations, the
free top and Poinsot's description; the Lagrange top.
\item Hamiltonian mechanics: the Legendre transform, Hamilton's equations on $T^*Q$, Poisson brackets,
canonical transformations and generating functions, the Hamilton--Jacobi equation.
\item Symmetry and reduction: momentum maps, Marsden--Weinstein reduction, the motion of a particle in a
central field revisited.
\item Integrable systems: the Liouville--Arnold theorem, action--angle variables; examples (the
two-body problem, the spherical pendulum, the Toda lattice); the KAM theorem (statement).
\item The $n$-body problem: central configurations, Sundman's theorem, the restricted three-body problem
and Lagrange points (survey).
\item Exercises with solutions.
\end{enumerate}
